\section{Method}
\label{sec:method}

\subsection{Overall Framework}
\begin{itemize}
    \item Follow the high-level path from native ICBHI cycles and SPRSound
    events through task alignment, shared encoding, hierarchical prediction,
    and separate native outputs.
    \item Separate shared representation learning from dataset-native
    evaluation.
\end{itemize}

\begin{figure*}[t]
    \centering
    \fbox{\parbox[c][1.20in][c]{0.96\textwidth}{
        \centering
        \textit{Figure 2 placeholder: native inputs $\rightarrow$ task
        alignment $\rightarrow$ shared BEATs representation $\rightarrow$
        eligibility-aware hierarchy $\rightarrow$ ICBHI flat-four and
        SPRSound binary outputs.}
    }}
    \caption{\textit{Figure 2 caption placeholder: high-level method flow;
    exclude optimizer details, auxiliary experiment branches, supporting
    datasets, and experiment identifiers.}}
    \label{fig:method_overview}
\end{figure*}

\subsection{Shared Acoustic Representation}
\begin{itemize}
    \item Introduce the pretrained BEATs encoder and its acoustic-tokenizer
    pretraining source.
    \item Describe the waveform-to-filterbank frontend, temporal aggregation,
    and shared projection used by both core tasks.
    \item State the downstream optimization scope of the encoder and shared
    projection.
    \item \textit{Equation placeholder: waveform-to-shared-representation
    mapping.}
    \item \textit{Citation placeholder: BEATs.}
\end{itemize}

\subsection{Hierarchical Prediction Nodes}
\begin{itemize}
    \item Define the Level-1 Normal/Abnormal probability.
    \item Define eligible Crackle and Wheeze attribute probabilities and their
    possible co-occurrence.
    \item Define the dependency between Level-1 and attribute predictions.
    \item \textit{Equation placeholder: Level-1 and attribute-node
    probabilities.}
\end{itemize}

\subsection{Eligibility-Aware Objective}
\begin{itemize}
    \item Define the node-availability indicator from the source annotation.
    \item Compute each node loss only on samples with an available target.
    \item Aggregate node losses without converting unavailable targets into
    negative labels.
    \item Aggregate ICBHI and SPRSound task losses under the fixed joint
    training contract.
    \item \textit{Equation placeholder: availability-masked node loss.}
    \item \textit{Equation placeholder: combined core classification
    objective.}
\end{itemize}

\subsection{Inherited Patient-Aware Objectives}
\begin{itemize}
    \item Introduce PAFA's patient cohesion-separation and global alignment
    objectives as inherited training components.
    \item Define their relationship to the shared representation and the
    training-only projection branch.
    \item State that the projection branch is removed at inference.
    \item \textit{Equation placeholder: inherited PAFA objective and total
    training objective.}
    \item \textit{Citation placeholder: PAFA.}
\end{itemize}

\subsection{Dataset-Native Readouts}
\begin{itemize}
    \item Reconstruct ICBHI Normal, Crackle, Wheeze, and Both predictions from
    the hierarchy using frozen validation-derived thresholds.
    \item Use the Level-1 output for SPRSound Normal/Adventitious prediction.
    \item Report the two native tasks with their own units and metrics.
    \item \textit{Equation placeholder: ICBHI flat-four decoding and
    SPRSound binary readout.}
\end{itemize}
